\section{Related Work}
\label{sec:related}

\paragraph{Object hallucination and contextual shortcuts.}
CHAIR measures unsupported object mentions in captions, while POPE evaluates
object existence with controlled negative-object sampling~\cite{chair,pope}.
Their complementary response formats make it possible to inspect both
generated descriptions and binary decisions. They do not, by themselves,
identify the computational origin of a failure. He
et al.~\cite{removedobjects} evaluate hallucinations after object removal and
propose object-focused preference optimization. Causal-HalBench formalizes
co-occurrence-induced spurious correlations and evaluates counterfactual
images~\cite{causalhalbench}. We build on these problem formulations; target
removal and contextual evaluation are not introduced here.

\paragraph{Activation steering in VLMs.}
VTI modifies visual and textual latent representations to reduce
hallucination~\cite{vti}. GrAInS derives steering directions from
gradient-based attribution of preferred versus dispreferred
outputs~\cite{grains}. REVIS constructs a purified visual direction and
calibrates sparse intervention, while RUDDER obtains an evidence direction
from prefill residual updates and applies adaptive
reinforcement~\cite{revis,rudder}. \method instead fits scalar gains on
selected, pre-projection attention-head outputs. This architectural
difference is an experimental choice, not evidence of superiority. It
requires comparison at matched preprocessing, decoding, calibration budget,
and object recall.

\paragraph{Localization and targeted parameter changes.}
HONES localizes task-relevant heads, attributes their downstream effects to
FFN neurons, and learns sparse neuron scaling~\cite{hones}. It is especially
close to our localization-then-scaling design. Seeing or
Knowing?~\cite{seeingknowing} examines visual-context sensitivity and connects
internal intervention analysis with steering. PURGE studies partition-aware
unlearning of spurious associations while preserving object
reasoning~\cite{purge}. Consequently, neither causal localization, selective
scaling, nor evidence--context separation is a first-of-its-kind claim for
\method. Its narrower candidate contribution is whether the
target-versus-control patch score selects better sites for bounded
head-output gains under information-preserving evaluation.

\paragraph{Scope of the comparison.}
Frozen-backbone steering and unlearning or preference optimization incur
different training costs and use different supervision. We report those
differences rather than treating all methods as equally resource-free.
Our four existing steering implementations are local adaptations, not
verified reproductions of published scores. Comparisons to HONES, the
visual-context steering method, and PURGE require additional implementation
or integration before a scientific claim can be made.
